
% =====================================================================
%  01  緒言（Introduction）  — 需要と供給の現状／反応式／設計の要点／生産目標
%  方針：略語は丸括弧で定義してから使う／"価値化"等の抽象語を避け具体化／
%        量とスペック（生産目標）を最下部の帯で明示／色は必要時のみ／横幅ぎりぎり。
%  内容は報告書 01_introduction.tex・02_process_overview.tex に準拠。
%  供給内訳 61/34/3% と PDH 収率 35-40% は 01_introduction.tex、
%  生産目標は 02_process_overview.tex Table（設計目標と原料仕様）。
% =====================================================================
\begin{frame}[t]

  \vspace*{3mm}   % 見出しを帯から離す（本文を scriptsize に縮小して縦を回収）
  % ===== 上：需要と供給の現状（1段落に圧縮・具体化）================
  {\bfseries\large プロピレンの需要と供給の現状}\par
  \vspace{0.3em}
  {\scriptsize\raggedright
    プロピレンはポリプロピレン・アクリロニトリル等の誘導品の基幹原料で、
    世界需要は年率約 $4$\,\% で拡大している。従来のスチームクラッキング・
    FCC（流動接触分解）の副産では伸びを賄えず、オンパーパス生産が
    世界能力の約 $2$ 割まで拡大している{\tiny\,（IHS Markit, Global Propylene Supply Demand, 2021, p.5）}。
    本設計は\textbf{PDH（プロパン脱水素）}によるオンパーパス生産を採用する。\par}

  \vspace{0.4em}

  % ===== 下：左=反応式（太枠）／右=設計の要点（課題 改行 →対策、4項目）====
  \begin{columns}[T,onlytextwidth]
    % ---- 左下：反応式（内容ぴったりの太枠で強調）-----------------
    \begin{column}{0.42\textwidth}
      {\bfseries\large 反応式}\par
      \vspace{0.35em}
      {\setlength{\fboxsep}{6pt}\setlength{\fboxrule}{1.6pt}%
       \fcolorbox{HeaderBlue}{white}{%
        \small\boldmath
        \renewcommand{\arraystretch}{1.08}%
        \begin{tabular}{@{}l@{}}
          主反応（脱水素化） \\[1pt]
          $\mathrm{C_3H_8 \rightleftharpoons {\color{SpRed}C_3H_6} + H_2}$ \\[5pt]
          副反応 \\[1pt]
          $\mathrm{C_3H_8 \rightarrow {\color{SpBlue}C_2H_4} + CH_4}$ \\[1pt]
          $\mathrm{{\color{SpBlue}C_2H_4} + H_2 \rightarrow {\color{SpBlue}C_2H_6}}$ \\[5pt]
          コーキング \\[1pt]
          $\mathrm{{\color{SpRed}C_3H_6} \rightarrow 3\,{\color{SpBlue}CH_{0.5}} + 2.25\,H_2}$
        \end{tabular}}}
    \end{column}
    % ---- 右下：設計の要点（課題 改行 →対策。対策は太字）----------
    \begin{column}{0.58\textwidth}
      {\bfseries\large 設計の要点}\par
      \vspace{0.35em}
      {\footnotesize\raggedright
        副生するオフガスの扱い\par
        $\to$\,\textbf{水素は {\color{SpRed}PSA} で回収し製品、軽質ガスは燃料}\par
        \vspace{0.28em}
        高温でコークが析出し触媒が分単位で失活\par
        $\to$\,\textbf{反応と再生を{\color{SpRed}周期的切替操作}で連続生産}\par
        \vspace{0.28em}
        プロピレンとプロパンの比揮発度が小さく分離困難\par
        $\to$\,\textbf{{\color{SpRed}膜分離}と蒸留のハイブリッドで分離}\par
        \vspace{0.28em}
        多数の設計変数が互いに依存し同時最適化が困難\par
        $\to$\,\textbf{全 23 変数を{\color{SpRed}ベイズ最適化}で同時決定}\par
      }
    \end{column}
  \end{columns}

  \vspace{0.4em}

  % ===== 最下部：生産目標の帯（量＋スペックを明示）================
  {\setlength{\fboxsep}{5pt}%
   \noindent\colorbox{HiBlue}{%
    \begin{minipage}{\dimexpr\linewidth-2\fboxsep\relax}
     \footnotesize\bfseries\color{HeaderBlue}%
     生産目標：原料 LPG \mbox{$\mathrm{C_3H_8}{:}\mathrm{C_4H_{10}}{=}9{:}1$} から、\\
     {\color{SpRed}プロピレン $99.5$ wt\%} を
     {\color{SpRed}年 $40$ 万 t（$1194$ kmol/h・稼働 $8000$ h/年）}製造
     {\scriptsize\mdseries\mbox{（世界の商業 PDH プラント 1 基 $65$–$90$ 万 t/年）}{\tiny\,（Borealis Kallo 2018・Fujian Meide 2022 ほか）}}
    \end{minipage}}}

\end{frame}
